\chapter{Les types de neurones selon leur fonction électrique}
\chaptermark{Les neurones comme composants}
\label{sec:eetypes}

Au chapitre~\ref{sec:neurontypes}, nous avons trié les neurones selon le neurotransmetteur que libère leur sortie. Un électronicien les trierait autrement: selon ce que le neurone fait de son signal, c'est-à-dire selon le composant qu'il réalise. Nous connaissons le composant principal du livre depuis le chapitre~\ref{sec:glutcomputer}: le neurone intégrateur à fuite~\eqref{eq:glutneuron}, un circuit RC suivi d'un comparateur. Mais le cerveau en contient d'autres. Le tableau~\ref{tab:eetypes} les réunit en quatre groupes, et presque tous ont déjà été rencontrés dans le livre.

\begin{center}
\footnotesize
\setlength{\tabcolsep}{3pt}
\renewcommand{\arraystretch}{1.15}
\begin{longtable}{>{\raggedright}p{3.1cm}>{\raggedright}p{3.3cm}>{\raggedright}p{4.7cm}>{\raggedright\arraybackslash}p{2.4cm}}
\caption{Les neurones comme composants: nom, équivalent en électronique, ce que fait le composant et où il apparaît dans le livre.}\label{tab:eetypes}\\
\toprule
Neurone & Composant & Ce qu'il fait & Où dans le livre \\
\midrule
\endfirsthead
\toprule
Neurone & Composant & Ce qu'il fait & Où dans le livre \\
\midrule
\endhead
\bottomrule
\endfoot
\multicolumn{4}{l}{\emph{Filtres et intégrateurs}} \\
neurone intégrateur à fuite & intégrateur RC avec comparateur & additionne les entrées sur une fenêtre $\tau$ et décharge au seuil & chap.~\ref{sec:glutcomputer}, \eqref{eq:glutneuron} \\
détecteur de coïncidences & ET avec fenêtre temporelle & ne décharge que si les entrées arrivent presque ensemble; $\tau$ très petit & chap.~\ref{sec:glutcomputer}, \ref{sec:code}, \eqref{eq:window} \\
neurone phasique & filtre passe-haut, détecteur de front & répond au changement de l'entrée, puis se tait & chap.~\ref{sec:sensors} \\
neurone à adaptation & filtre passe-haut avec contrôle automatique de gain & la fréquence baisse sous une entrée constante & chap.~\ref{sec:sensors}, \eqref{eq:nakarushton} \\
résonateur & circuit oscillant, filtre passe-bande & répond le mieux à une entrée de sa propre fréquence & section~\ref{sec:resonator} \\
intégrateur sans fuite & intégrateur à amplificateur opérationnel & transforme une vitesse en position et la maintient & section~\ref{sec:perfectint} \\
\midrule
\multicolumn{4}{l}{\emph{Convertisseurs}} \\
neurone tonique & convertisseur tension-fréquence & la fréquence suit l'entrée linéairement et se fatigue à peine & chap.~\ref{sec:code} \\
neurone gradué (sans impulsions) & amplificateur analogique, tampon & transmet une tension continue sur une courte distance & chap.~\ref{sec:sensors} \\
neurone redresseur & redresseur simple alternance & la fréquence n'est jamais négative, $[u]_+$ & chap.~\ref{sec:glutcomputer}, \ref{sec:gaba} \\
paire \og ON-OFF\fg{} & paire de redresseurs, code à deux fils & l'un répond à \og plus\fg{}, l'autre à \og moins\fg{} & chap.~\ref{sec:glutcomputer}, \eqref{eq:dualrail} \\
neurone à inhibition shuntante & diviseur analogique, réglage du gain & divise l'entrée au lieu de la soustraire & chap.~\ref{sec:gaba}, \eqref{eq:shunt} \\
\midrule
\multicolumn{4}{l}{\emph{Générateurs et temps}} \\
stimulateur & oscillateur à relaxation, multivibrateur & émet de lui-même des impulsions à fréquence constante & chap.~\ref{sec:serocomputer}, \ref{sec:dofa} \\
stimulateur à entrée & oscillateur commandé en tension & rythme propre que l'entrée accélère ou ralentit & chap.~\ref{sec:chemoutput} \\
neurone à bouffées & générateur de bouffées déclenché & émet des bouffées d'impulsions rapprochées, les \og traits\fg{} & chap.~\ref{sec:code}, \eqref{eq:burst} \\
neurone calé sur la phase & détecteur de phase, boucle à verrouillage de phase & décharge à une phase donnée de l'onde d'entrée & chap.~\ref{sec:sensors}, \ref{sec:chemoutput} \\
axone à retard & ligne à retard & retarde le signal d'un temps donné & chap.~\ref{sec:glutcomputer}, fig.~\ref{fig:itd} \\
\midrule
\multicolumn{4}{l}{\emph{Mémoire et commutation}} \\
neurone bistable & trigger de Schmitt, verrou & une brève entrée le fait passer dans un état durable & section~\ref{sec:bistable} \\
neurone à branches dendritiques & multiplexeur, petit réseau multicouche & une branche ne laisse passer l'entrée que si une entrée d'autorisation est présente & chap.~\ref{sec:synthesis} \\
\end{longtable}
\end{center}

Une réserve compte plus que tout le tableau: il ne s'agit pas de cellules différentes, mais de \emph{modes} différents. Un même neurone peut être un intégrateur quand l'entrée est lente et un détecteur de coïncidences quand l'inhibition a raccourci son $\tau$ (chapitre~\ref{sec:code}); un stimulateur pendant la veille et un récepteur silencieux pendant le sommeil. Le composant obtenu dépend des canaux ioniques de la membrane et des canaux de diffusion qui les ajustent.

\section{Quatre composants que le livre n'a pas encore rencontrés}

\subsection{Le résonateur}
\label{sec:resonator}

Certains neurones possèdent un courant lent qui s'oppose à toute variation de tension: si la tension monte, le courant la tire vers le bas, si elle baisse, il la tire vers le haut, mais avec retard. Pour un électronicien, un tel courant se comporte comme une inductance, et avec la capacité de la membrane il forme un circuit oscillant. Le neurone répond le plus fortement à une entrée d'une fréquence donnée, en général de quelques hertz à une vingtaine, et plus faiblement aux entrées plus lentes ou plus rapides~\cite{HutcheonYarom2000}. C'est un filtre passe-bande dans une seule cellule: dans une entrée bruitée, il sélectionne le rythme de sa fréquence, par exemple le rythme local du chapitre~\ref{sec:gaba}.

\subsection{L'intégrateur sans fuite}
\label{sec:perfectint}

Un intégrateur à fuite ne se souvient de son entrée que pendant un temps $\tau$, quelques dizaines de millisecondes. Mais pour rester immobile, l'œil doit se souvenir de sa position pendant des secondes: l'ordre de \og tourner\fg{} arrive comme une vitesse brève, alors qu'il faut ensuite tenir l'œil dans la nouvelle position. Le réseau résout ce problème par la même rétroaction que la mémoire du chapitre~\ref{sec:glutcomputer}. Si un groupe de neurones de constante de temps $\tau$ s'excite lui-même avec un poids $w$, alors $\tau\,\dd r/\dd t=-r+w\,r+I$, et la constante de temps du groupe vaut
\begin{empheq}[box=\keybox]{equation}
\tau_{\text{eff}} \;=\; \frac{\tau}{1-w} .
\label{eq:perfectint}
\end{empheq}
Avec $\tau=0.1\,\text{s}$ et $w=0.995$, on obtient $20\,\text{s}$: un intégrateur presque idéal fait de neurones dont chacun, seul, oublie en un dixième de seconde~\cite{Seung1996}. Le prix à payer est un réglage précis: pour $w$ un peu supérieur à un, l'intégrateur part en saturation; un peu inférieur, il oublie vite. Un tel intégrateur a donc besoin d'un réglage permanent par apprentissage, comme celui que décrit le chapitre~\ref{sec:motorlearning}.

\subsection{Le neurone bistable}
\label{sec:bistable}

Aux chapitres~\ref{sec:glutcomputer} et~\ref{sec:gaba}, la bascule était assemblée à partir d'un groupe de neurones. Mais il existe aussi des bascules dans une seule cellule. Certains neurones possèdent un courant qui, une fois enclenché, entretient lui-même l'excitation: une brève entrée fait passer le neurone dans une activité durable, un \emph{plateau}, et l'inhibition le ramène au repos. C'est un trigger de Schmitt: le neurone a deux états stables et un hystérésis entre eux. C'est le cas, par exemple, des neurones \nt{ACET} qui commandent les muscles, et cette propriété est activée par un canal de diffusion: le plateau apparaît quand la sérotonine atteint le neurone~\cite{Hounsgaard1988}. La même cellule est un verrou avec \nt{SERO} et un intégrateur ordinaire sans lui.

\subsection{Intégrateurs et résonateurs}

La théorie répartit les cellules excitables en deux classes~\cite{Izhikevich2007}. Les \emph{intégrateurs} ressemblent à un oscillateur commandé en tension qui part de zéro: juste au-dessus du seuil, le neurone décharge aussi lentement qu'on veut, et sa fréquence croît continûment avec l'entrée. Un tel neurone additionne tout ce qui arrive et transmet bien l'amplitude de l'entrée par sa fréquence. Les \emph{résonateurs} ressemblent à un oscillateur à fréquence minimale: ils ne savent pas décharger lentement, émettent d'emblée des impulsions à fréquence finie dès que l'entrée atteint le seuil, et préfèrent une entrée de leur propre fréquence. Les premiers sont les composants naturels du code de fréquence du chapitre~\ref{sec:code}, les seconds du code temporel.

\begin{keyblock}
\begin{proposition}[Un neurone, de nombreux composants]
\label{prop:eetypes}
Selon leur fonction électrique, les neurones fonctionnent comme des intégrateurs à fuite ou sans fuite, des détecteurs de coïncidences, des filtres passe-haut et passe-bande, des convertisseurs tension-fréquence, des redresseurs, des diviseurs, des oscillateurs, des lignes à retard et des bascules. Le composant que donne une cellule dépend de ses canaux ioniques et des canaux de diffusion qui les ajustent. Les modes eux-mêmes sont mesurés; la mesure dans laquelle le cerveau exploite cette reconfiguration pour calculer est surtout une hypothèse.
\end{proposition}
\end{keyblock}

Pour un ingénieur, cela ressemble à un réseau logique programmable: le matériel est le même, et c'est la configuration qui fixe le composant. Sauf qu'ici on ne change pas la configuration au moment de la programmation, mais en marche, par les canaux de diffusion lents dont il a été question aux chapitres~\ref{sec:serocomputer}--\ref{sec:acet}.

\section{Exercices}

\begin{exercise}[\lvlA]
\label{ex:18:1}
\emph{Tolérance de l'intégrateur sans fuite.} Des neurones avec $\tau=0.1$~s tiennent la position de l'œil selon~\eqref{eq:perfectint} pendant $T=1$~s avec une erreur d'au plus $\pm5\%$. (a)~Trouvez la plage admissible de $w$. (b)~$w$ est la somme des poids de $K$ synapses, chacun avec une erreur indépendante de $10\%$. Pour quel $K$ la dispersion de la somme tient-elle dans la tolérance?
\end{exercise}

\begin{exercise}[\lvlB]
\label{ex:18:2}
\emph{Le résonateur comme circuit.} Décrivons le courant lent de la section~\ref{sec:resonator} par une variable $W$: $C\,\dd V/\dd t=-g_LV-g_wW+I$, $\tau_w\,\dd W/\dd t=V-W$. Montrez que c'est un circuit $RC$ avec une branche parallèle $R_w=1/g_w$, $L_w=\tau_w/g_w$; pour $C/g_L=10\ms$, $\tau_w=100\ms$, $\alpha=g_w/g_L=4$, trouvez la fréquence de résonance et le rapport $|Z|/|Z(0)|$ en ce point.
\end{exercise}

\begin{exercise}[\lvlB]
\label{ex:18:3}
\emph{Comment un intégrateur se met à décharger.} (a)~Neurone $\tau\,\dd V/\dd t=-V+\mu$, $\tau=10\ms$, de seuil $\theta$ avec remise à zéro. Pour quel $\mu/\theta-1$ la fréquence vaut-elle $1$~Hz? (b)~Quelle fréquence donne le modèle $\tau\,\dd v/\dd t=v^2+\mu$ (impulsion: départ de $v$ vers $+\infty$, remise à $-\infty$) pour $\mu=0.01$?
\end{exercise}

\begin{exercise}[\lvlB]
\label{ex:18:4}
\emph{Simulation: intégrateur et résonateur.} Le modèle de l'exercice~\ref{ex:18:2} avec $g_L=1$, $\tau_m=10\ms$, $\tau_w=100\ms$, un seuil $1$ et une remise $V\to0$ ($W$ inchangé) reçoit $\mu=1.5\sin2\pi ft$, $f=1$--$40$~Hz. Dans quelle bande de fréquences le neurone émet-il des impulsions pour $\alpha=0$ et pour $\alpha=4$? Comparez avec la condition $1.5\,|Z(f)|>1$.
\end{exercise}

\begin{exercise}[\lvlB]
\label{ex:18:5}
\emph{Un verrou fait d'une seule cellule.} Neurone à courant de plateau: $\tau\,\dd r/\dd t=-r+F(wr+I)$, $F(x)=\min(1,\max(0,x))$, $\tau=10\ms$, $w=1.5$, fond $I=-0.2$. (a)~Quelle durée minimale d'une impulsion $I=0.3$ fait passer le neurone de $r=0$ au plateau? (b)~Quelle amplitude minimale $B$ d'une longue impulsion $I=-0.2-B$ le ramène au repos?
\end{exercise}

\begin{exercise}[\lvlB]
\label{ex:18:6}
\emph{Autoréglage de l'intégrateur.} Pendant la fixation du regard, l'entrée est $I=0$, un capteur de glissement donne la vitesse de dérive $e=\dd r/\dd t$, et $\dd w/\dd t=-\eta\,e\,r$. Pour $w=0.95$, $\tau=0.1$~s, $\langle r^2\rangle=1$, trouvez $\eta$ tel qu'après $60$~s de fixations $|1-w|$ entre dans la tolérance de l'exercice~\ref{ex:18:1}. Que se passe-t-il si l'on apprend aussi pendant les rotations de l'œil?
\end{exercise}

\begin{exercise}[\lvlC]
\label{ex:18:7}
\emph{Pourquoi reconfigurer le composant?} Un canal de diffusion commute l'$\alpha$ de l'exercice~\ref{ex:18:2} entre $0$ et $4$. De quel facteur le résonateur l'emporte-t-il sur l'intégrateur en rapport signal/bruit pour une faible sinusoïde dans un bruit blanc de courant à $11$~Hz et à $1$~Hz? Imaginez une tâche où la commutation de mode elle-même est avantageuse.
\end{exercise}
